% Working draft for a future Musings post.
% Compile from the repository root with:
%   mkdir -p /tmp/visions-of-the-future-of-science
%   pdflatex -interaction=nonstopmode -output-directory /tmp/visions-of-the-future-of-science drafts/visions-of-the-future-of-science.tex
%
% Once the text is ready, adapt it to the site's Markdown front matter and
% publish it under _posts/ with `section: musings`.
\documentclass[11pt,a4paper]{article}

\usepackage[T1]{fontenc}
\usepackage[utf8]{inputenc}
\usepackage[english]{babel}
\usepackage{lmodern}
\usepackage{microtype}
\usepackage[a4paper,margin=30mm]{geometry}
\usepackage{hyperref}
\usepackage{xcolor}

% Assistant-proposed prose is deliberately visible while this remains a
% collaborative working draft.  Remove \AI{} when the wording is accepted.
\newcommand{\AI}[1]{{\color{blue}\textit{AI draft: #1}}}

\hypersetup{
  colorlinks=true,
  linkcolor=black,
  urlcolor=blue,
  pdftitle={Visions of the Future of Science},
  pdfauthor={Leo Shaposhnik}
}

\title{Visions of the Future of Science}
\author{Leo Shaposhnik}
\date{\today}

\begin{document}

\maketitle

% This file deliberately starts as a clean writing scaffold.  Draft prose below
% can be copied into the eventual Markdown post with only light formatting.

\begin{abstract}
What should science become when AI systems can participate in the work of
discovery?  This essay considers the opportunities, responsibilities, and
institutional choices that will shape that answer.
\end{abstract}

\section{The question}

\subsection{From papers to an integrated body of knowledge}

I would like to argue for the emergence of a new type of knowledge
representation away from papers towards something more integrated that we can
parse, inspect, and add to in a modular fashion.

\AI{The question is not merely what AI will allow individual scientists to do
faster, but what kind of scientific record we should build for a world in which
both people and agents can read, compare, and extend an enormous body of work.
At present, scientific knowledge is primarily preserved as papers: texts first
shared as preprints and, often later, placed in journals after peer review.
Papers are connected mainly by citations.  Yet a citation usually points to an
entire paper rather than to the particular definition, calculation, argument,
or qualification on which later work relies.  The resulting literature is
extraordinarily valuable, but it is only loosely interconnected and difficult
to traverse at the level where scientific arguments actually depend on one
another.}

\AI{This is a suboptimal representation of knowledge, especially now that
agents can work through large amounts of text quickly.  It is also an
opportunity.  Rather than asking machines merely to summarize or reproduce the
existing literature, we could use them to help construct a representation of
scientific knowledge that is more explicit about what is assumed, what is
claimed, and how one claim is transformed into another.  Such a representation
would be useful not only for scale, but also for rigor: it would make the
dependencies of an argument easier to inspect and its weakest assumptions
easier to locate.}

\AI{The closest existing model is a proof assistant such as Lean.  In a mature
formal theory, typed statements and verified proofs make relations between
results mechanically explicit.  This is an important ideal, but it is not yet
a universal template for theoretical physics.  Many physical theories---string
theory is one example---are not specified sharply enough for their full content
to be formalized.  Moreover, formalization is not merely the act of writing
code: one must also judge whether the formal object faithfully represents the
informal idea.  The aim need therefore not be that every scientific claim
immediately compiles.}

\AI{A more attainable intermediate goal is a typed or semi-typed account of
scientific reasoning.  A theory contains assumptions, claims, and procedures
that transform some claims into others.  Those procedures can have different
epistemic status: a heuristic argument, a rigorous derivation, a numerical
calculation, a symbolic computation, or an empirical inference.  In this
limited sense, scientific work resembles a network of morphisms: claims are
related by transformations whose inputs, outputs, and status should be made
visible.  A new result could then be integrated not just by attaching another
document to the archive, but by recording how it depends on and modifies the
structure already present.}

\AI{The representation should be developed at a scale where shared standards
are realistic.  A research community could define the objects, claims, and
argument forms that matter in its own field: quantum circuits and inequalities
in quantum information, for example.  The point is not to impose a single
universal ontology on all science.  It is to make local scientific languages
explicit enough that work becomes modular.  Individual projects can contribute
small, reviewable pieces to a shared whole, while the relations between those
pieces remain available for people and machines to follow.}

\AI{Lamport-style structured arguments suggest one practical model.  Each step
of an argument can be presented either as an introduced assumption or as a
claim justified by specified earlier steps.  Even when the justification is not
a formally checked proof, this exposes the flow of reasoning.  A representation
built from such steps would allow a reader---or an agent---to traverse an
argument, inspect individual dependencies, and see precisely where trust is
being placed.}

\AI{The largest gain would be a more reliable foundation for agent-assisted
research.  As long as a theory is not fully formalized, our confidence rests on
the assumptions we accept and the transformations of those assumptions that we
trust.  We should therefore make both as inspectable as possible.  A shared
knowledge base with explicit provenance and dependencies would give agents a
body of work they can cite, query, and extend in ways that are easier to verify
than a free-standing generated text.  It would also let scientists check the
relevant parts of an argument without having to reconstruct the whole history
of a field.}

\AI{This vision has institutional consequences.  If a contribution is a
reviewable extension of an open, living knowledge artifact, then scientific
credit and review need not be organized solely around the publication of a
closed paper in a journal.  Review could become a visible collective practice
of maintaining the reliability of shared infrastructure.  Such an approach
would support open science, concurrent collaboration, and larger projects in
which different groups contribute interoperable pieces.  The challenge is to
design the type systems, governance, incentives, and interfaces that make this
possible without turning scientific judgment into a bureaucratic schema.}

\section{What becomes abundant}

% Discuss tasks that AI can make cheaper or faster, without equating abundance
% of output with abundance of understanding.

\section{What remains scarce}

% Discuss judgment, taste, trust, verification, problem choice, and collaboration.

\section{Institutions for a better scientific future}

% Develop concrete visions for incentives, shared infrastructure, and evaluation.

\section{An invitation}

% Close with the constructive question or proposal you want readers to take away.

\end{document}
